
@@M0@@ la classe d'un tel faisceau. On a alors :
@@M1@@
Comme @@M2@@, on obtient, d'après la définition de @@M3@@ :
@@M4@@

Par ailleurs, @@M5@@ est d'augmentation @@M6@@; la classe de Chern qui intervient dans la formule 4.4 relative à @@M7@@ est donc @@M8@@. Compte tenu de la définition de @@M9@@, et de 3.5, la relation à démontrer s'écrit :
@@M10@@
et est donc vérifiée.

\textbf{Remarque 4.4.1.} Il est possible aussi de déduire 4.4 d'un énoncé purement algébrique correspondant, cf. @@M11@@ I 1.3.

\textbf{4.5.} Nous allons montrer maintenant que pour montrer le théorème de Riemann--Roch dans le cas envisagé dans ce paragraphe, on peut se ramener au cas où @@M12@@ est un diviseur de @@M13@@, et où l'élément auquel s'applique la formule est image inverse d'un élément de @@M14@@.

D'après 4.2, la formule à démontrer :
@@M15@@
est équivalente à la formule :
@@M16@@
Compte tenu de la fonctorialité de @@M17@@, et de 4.3, elle peut encore s'écrire :
@@M18@@
soit, d'après VII 3.4 :
@@M19@@

